\begin{abstract}
Parallel grippers are reliable and inexpensive, but once they pinch an object they cannot change its pose within the grasp. Rather than adding sensing, finger actuation, or contact-rich planning, we add passive geometry to the robot's workspace, and we show that this geometry can be designed in closed form. The design starts from what passive contacts can and cannot do: they cannot determine a motion that both slides and turns the object in the grip, so a planar change of grasp is executed as a shift, during which the tool holds the object while the gripper slides along it, followed by a turn, during which the gripper orbits one contact while the grasp holds. Two rods spinning freely on bearings provide both motions. Two rods are the fewest that hold an arbitrary shift; one of them then guides the turn, and the other separates by itself. The rod placement is a projection at one vertex of a convex or nonconvex polygonal outline, the gripper path is a segment followed by a circular arc, and a single quadratic equation tells whether an installed pair can serve a request. No friction coefficient or grip-force value enters the design, no sensing is used during the manipulation motion, and because the rods spin freely, grasp slip during the turn is predicted from the pad size and the lever arm alone. Hardware turns on a Kinova Jaco reached mean absolute errors of $3.6^{\circ}$--$7.5^{\circ}$; in simulation, complete sequences ended within $1.6^{\circ}$ and 0.44 mm of the request, and the measured tilt, slip, and separation agreed with the closed-form predictions.
\end{abstract}
